\documentclass[10pt,a4paper]{amsart}
\usepackage[margin=19mm]{geometry}
\usepackage{amsmath,amssymb,mathtools}
\usepackage[scr=rsfs,cal=euler]{mathalfa}
\usepackage{xfrac}
\usepackage[unicode,hidelinks,bookmarks=false]{hyperref}
\newcommand{\ie}{i.e.\ }
\newcommand{\cf}{cf.\ }
\newcommand{\resp}{respectively\ }
\newcommand{\Subsection}{\subsection}
\providecommand{\nobreakdash}{\nobreak\hbox{-}}
\setlength{\emergencystretch}{2em}
\multlinegap0pt
\usepackage{amssymb}
\newtheorem{lem}{Lemma}
\newcommand{\glob}[1]{\ensuremath{\langle\, #1\,\rangle}}
\newcommand{\st}{\mbox{\ensuremath{\ast}}}
\newcommand{\sym}{\mathfrak{S}}
\title{Descent generating polynomials for ($n-3$)- and ($n-4$)-stack-sortable (pattern-avoiding) permutations}
\author{Sergey Kitaev, Philip B. Zhang}
\date{Source: arXiv:2503.22067v3 (CC BY 4.0)}
\begin{document}
\maketitle

\noindent\emph{Source and licence.} Descent generating polynomials for ($n-3$)- and ($n-4$)-stack-sortable (pattern-avoiding) permutations. Version arXiv:2503.22067v3 (CC BY 4.0), distributed by arXiv under the Creative Commons Attribution 4.0 International licence, http://creativecommons.org/licenses/by/4.0/, as recorded in the arXiv licence metadata for this exact version and shown on the abstract page https://arxiv.org/abs/2503.22067v3. Full licence notice: \texttt{licenses/arxiv-2503.22067-CC-BY-4.0.txt}.\par\smallskip

\noindent\textit{Editorial scope.} Counted unit: the article's lem lem-length-3 (source0.tex lines 608--614). The article's complete original proof of it is delivered with it (source0.tex lines 615--666, one \texttt{proof} environment of 52 lines). No mathematics is written, repaired or re-derived: the statement and the proof are the article's own lines, in the article's own notation, and no retained line is substituted or re-flowed.  Retained uncounted as the source-local context the proof needs: lines 376--380.  Nothing was omitted from any retained span, so the package carries no omission record.  No retained line contains a citation, so the wrapper carries no bibliography and no bibliography entry.  No retained line contains a figure environment, an image-inclusion command or drawing code, so no image is delivered and the package declares no asset dependency.  Editorial formatting only, in addition: an AMS \texttt{amsart} preamble that restates the article's own declarations and macros used by the retained lines, namely the environments \texttt{lem} on separate counters, together with the standard packages those lines need.  The retained environments print in this document as Lemma 1 in order of appearance, whereas the article prints the same items with the numbering of its own full document.  The complete native source of the article as distributed in the arXiv e-print for this exact version is bundled unchanged as \texttt{sources/175535/source0.tex}.
\begin{lem}\label{lem-1}
	For any $n\ge 2$, we have  that
	$$A_n \glob{\st n\st (n-1)\st }
		= \frac{1}{2}\big( A_{n}(x) + (x-1) A_{n-1}(x) \big).$$
\end{lem}

\begin{lem}\label{lem-length-3}
	For all $n\ge 3$,
	\begin{align*}
		A_n \glob{\st n\st (n-1)\st (n-2) \st }
		= \frac{1}{6} \left( A_{n}(x) + 3 (x-1) A_{n-1}(x) + 2(x-1)^2 A_{n-2}(x) \right).
	\end{align*}
\end{lem}

\begin{proof}
	We note that the set of all permutations of $[n]$ is the disjoint union of the following three types, where $abc$ forms a permutation of $\{n-2,n-1,n\}$:
	
	(i) $\glob{\st\,a\,?\st\,b\,?\st\,c\,\st}$;
	
	(ii) $\glob{\st\,a\,b\,?\st\,c\,\st}$ or $\glob{\st\,a\,?\st\,b\,c\,\st}$;
	
	(iii) $\glob{\st\,a\,b\,c\,\st}$.
	
	The strict-gap form of the computation in Lemma~\ref{lem-1} is
	\[
	A_{n-1}\glob{\st(n-1)\,?\st(n-2)\st}
	=\frac12\left(A_{n-1}(x)-(x+1)A_{n-2}(x)\right).
	\]
	This polynomial is unchanged if the two distinguished top letters are interchanged, because the enforced gap makes all their adjacent neighbours smaller.
	Thus
	\begin{align*}
		A(\sym_n\cap\glob{\st a\,b\,?\st\,c\,\st})
		&=x^{\chi(a>b)}\frac12\left(A_{n-1}(x)-(x+1)A_{n-2}(x)\right),\\
		A(\sym_n\cap\glob{\st a\,?\st\,b\,c\,\st})
		&=x^{\chi(b>c)}\frac12\left(A_{n-1}(x)-(x+1)A_{n-2}(x)\right),\\
		A(\sym_n\cap\glob{\st a\,b\,c\,\st})
		&=x^{\chi(a>b)+\chi(b>c)}A_{n-2}(x),
	\end{align*}
	where $\chi(\cdot)$ is $1$ if the statement is true and $0$ otherwise. Since relabelling the three separated distinguished letters in type (i) does not change descents, the six orders have the same generating polynomial. Hence
	\begin{align*}
		&6A(\sym_n\cap\glob{\st a\,?\st\,b\,?\st\,c\,\st})
		+6(x+1)\frac12\left(A_{n-1}(x)-(x+1)A_{n-2}(x)\right)\\
		&\qquad +(x^2+4x+1)A_{n-2}(x)=A_n(x),
	\end{align*}
	and therefore
	\begin{align*}
		A_n\glob{\st a\,?\st\,b\,?\st\,c\,\st}
		=\frac16\left(A_n(x)-3(x+1)A_{n-1}(x)+2(x^2+x+1)A_{n-2}(x)\right).
	\end{align*}
	
	The set counted in the lemma is the disjoint union
	\begin{align*}
		&\glob{\st n\,?\st(n-1)\,?\st(n-2)\st}
		\mathbin{\bigcup}\glob{\st n(n-1)\,?\st(n-2)\st}\\
		&\qquad\mathbin{\bigcup}\glob{\st n\,?\st(n-1)(n-2)\st}
		\mathbin{\bigcup}\glob{\st n(n-1)(n-2)\st}.
	\end{align*}
	Consequently,
	\begin{align*}
		A_n\glob{\st n\st(n-1)\st(n-2)\st}
		={}&\frac16\left(A_n(x)-3(x+1)A_{n-1}(x)+2(x^2+x+1)A_{n-2}(x)\right)\\
		&+2\frac{x}{2}\left(A_{n-1}(x)-(x+1)A_{n-2}(x)\right)+x^2A_{n-2}(x)\\
		={}&\frac16\left(A_n(x)+3(x-1)A_{n-1}(x)+2(x-1)^2A_{n-2}(x)\right).
	\end{align*}
	This completes the proof.
\end{proof}

\end{document}
